\section*{Abstract}

Deep learning training on high-performance computing (HPC) clusters is conventionally
treated as a problem of GPU throughput, and optimisation effort is directed
at the model.
For semantic segmentation over remotely sensed imagery, however, the training corpus
is typically stored as very large numbers of small image and mask files on a shared
parallel filesystem, and the cost of reaching those files is borne once per sample.
%
We show that for a representative archaeological site-detection workload
- the detection of mounded settlement sites (\emph{tells}) in the southern
Mesopotamian floodplain~\cite{Casini23} -
the GPU is frequently idle while the data path is the binding constraint:
across three architectures, time spent waiting on data exceeds forward-and-backward
compute time by 8--20$\times$ under a default loader configuration, and two loader
flags alone move GPU utilisation by 24--51 percentage points. Segmentation quality
moves far less under the same change, and not consistently in one direction, which
is itself informative: this workload's headroom is in throughput, not accuracy.
%
We introduce \pads{} (Profiler-Adaptive Data Staging), a drop-in replacement for the
PyTorch \code{DataLoader} that profiles its own behaviour over a short warm-up window,
classifies the dominant bottleneck from the measured timings, selects a staging policy
- loose files, tar shards, or node-local scratch staging - and executes that policy
with a prefetch depth derived from the measurements.
%
Because the sample permutation is fixed at the start of each epoch, the identity of
future file accesses is known exactly, so prefetching is exact rather than speculative.
%
We evaluate \pads{} against an unmodified loader, a manually tuned loader, fixed
tar-shard streaming, and a full pre-staging upper bound, on three segmentation
architectures spanning a CNN encoder--decoder, an attention-augmented CNN, and a
transformer.
The staging mechanism itself delivers on that promise when actually invoked: tar
shards alone beat an unmodified loader on GPU utilisation for every architecture
tested, and a full pre-staging ceiling reaches the highest GPU utilisation of any
configuration tested for two of three architectures. \pads{}'s adaptive policy selection, however, activated staging in only a small minority of the configurations where the fixed-policy experiments showed that staging could provide a measurable benefit. This behaviour was traced to a specific and quantified defect in the profiler's warm-up measurement. One-time GPU initialisation costs dominated the short averaging window, inflating the measured GPU compute time and systematically understating the true data-to-compute time ratio by approximately an order of magnitude. The resulting conservative policy selection therefore reflects a limitation of the profiling procedure rather than a failure of the staging mechanism itself.
. A validated correction, discarding warm-up
batches before measuring, recovers full-epoch measurement accuracy at negligible
added cost.
%
Alongside the positive throughput claim we make a negative accuracy claim: a pipeline
optimisation is only sound if segmentation quality is preserved, and we therefore
report intersection-over-union against a reference reproduction throughout.

\vspace{1em}
\textbf{Keywords:} Data Staging; I/O Bottlenecks; Semantic Segmentation; Remote Sensing;
High-Performance Computing
